% Audit complet de la branche de mise à niveau : mod_ep, mod_epsynthesis, mod_syllabus,
% mod_syllabusstaff, mod_syllabusstudent.
% Compilation : latexmk -pdf audit-complet.tex
\documentclass[11pt,a4paper]{article}
\usepackage[french]{babel}
\usepackage[T1]{fontenc}
\usepackage[utf8]{inputenc}
\usepackage{lmodern,microtype,booktabs,tabularx,xltabular,xcolor,hyperref,array,enumitem}
\usepackage[margin=2.2cm]{geometry}
\hypersetup{colorlinks=true, linkcolor=blue!50!black, urlcolor=blue!50!black}
\newcommand{\code}[1]{\texttt{\detokenize{#1}}}
\newcommand{\grav}[1]{\textsc{#1}}
\setlist{itemsep=2pt, topsep=4pt}
\title{Audit complet de la mise à niveau d'EP Manager et des modules syllabus}
\author{Sébastien Lefebvre}
\date{Novembre 2026 -- mod\_ep 1.0.6, mod\_epsynthesis 1.0.3, modules syllabus 1.5.1}

\begin{document}
\maketitle

\section{Objet et méthode}

L'audit porte sur l'ensemble de la branche de mise à niveau, avant sa fusion : les cinq modules
d'activité (\code{mod_ep}, \code{mod_epsynthesis}, \code{mod_syllabus},
\code{mod_syllabusstaff}, \code{mod_syllabusstudent}, environ 45\,000 lignes de PHP), leur
liaison avec \code{mod_stage}, leurs tests, la CI et la documentation. Il complète les rapports
de charge, de concurrence et de bout en bout déjà présents dans \code{docs/audit}, sans en
reprendre les mesures.

\paragraph{Six revues du code}, menées en parallèle, chacune sur tout le périmètre de son
domaine, en lecture seule, chaque constat tracé dans le code de bout en bout :
\begin{enumerate}
  \item sécurité de \code{mod_ep} et \code{mod_epsynthesis} (droits par page et par action,
    identifiants de la requête, clé de session, échappement, fichiers, imports) ;
  \item sécurité des modules syllabus (mêmes points, plus l'anonymat des avis et les actions
    du suivi sur les modules d'autres cours) ;
  \item intégrité des données (schéma et mises à jour, suppression, réinitialisation,
    sauvegarde et restauration, RGPD, concurrence, tâches planifiées) ;
  \item performance et portabilité SQL (PostgreSQL, MariaDB/MySQL, SQL Server, Oracle) ;
  \item langues, accessibilité et cohérence de l'interface ;
  \item tests, CI, métadonnées de version et cohérence de la documentation.
\end{enumerate}
Chaque constat retenu a ensuite été revérifié dans le code ; les défauts confirmés et de portée
limitée sont corrigés, avec un test qui échoue sans la correction quand c'est possible. Les
autres sont décrits en section~\ref{sec:residuel}, avec la raison pour laquelle ils restent.

\paragraph{Contrôles sur le banc} (Moodle 4.5, PHP~8.3, PostgreSQL~16, jeu de données de 1\,000
étudiants) : schéma de la base comparé aux fichiers \code{install.xml} après les mises à jour
successives, test de conformité RGPD du cœur, balayage dynamique des droits
(\code{docs/audit/banc/balayage_droits.py}, nouveau), scénario de bout en bout des modules
syllabus rejoué avec le code corrigé.

\section{Synthèse}

\begin{center}
\begin{tabularx}{\linewidth}{@{}X r r r@{}}
\toprule
\textbf{Domaine} & \textbf{Constats} & \textbf{Corrigés} & \textbf{Restants} \\
\midrule
Sécurité & 14 & 8 & 6 \\
Intégrité des données & 15 & 9 & 6 \\
Performance et portabilité & 7 & 5 & 2 \\
Langues et accessibilité & 8 & 4 & 4 \\
Tests, CI et documentation & 7 & 2 & 5 \\
\midrule
\textbf{Total} & \textbf{51} & \textbf{28} & \textbf{23} \\
\bottomrule
\end{tabularx}
\end{center}
Un constat peut regrouper plusieurs points de même nature (par exemple trois écarts de
portabilité SQL, ou les tableaux sans conteneur défilant de cinq pages).

Aucun constat critique. Les plus graves, tous corrigés :
\begin{itemize}
  \item un compte qui voit tous les dossiers d'une activité EP (la DEVE) pouvait lire le
    supplément au diplôme -- nom, numéro, date de naissance -- de n'importe quel compte du
    site, en changeant l'identifiant dans l'adresse ; reproduit sur le banc avant correction ;
  \item un étudiant pouvait faire sauter l'étape DEVE d'une déclaration en changeant son type ;
  \item un responsable pouvait lire les avis avant la fin du recueil (fermeture puis
    réouverture de la fenêtre), et l'heure de chaque avis, dans le journal, permettait de
    retrouver son auteur en comparant les résultats ;
  \item une restauration rendait invérifiables les suppléments déjà remis aux étudiants.
\end{itemize}
Les constats restants sont de gravité moyenne ou basse ; deux recommandations de gravité
haute concernent la CI (mise à jour depuis la version en production, scénarios Behat de
\code{mod_ep}).

\section{Contrôles automatiques}

\begin{center}
\begin{tabularx}{\linewidth}{@{}X l@{}}
\toprule
\textbf{Contrôle} & \textbf{Résultat} \\
\midrule
Schéma de la base du banc, mis à jour version après version, comparé aux \code{install.xml}
(\code{check_database_schema.php}) & conforme \\
Test de conformité RGPD du cœur (fournisseurs des cinq modules compris) & 1\,648 tests réussis \\
Tests unitaires des cinq modules (132 + 8 + 43 + 6 + 1), après corrections & 190 réussis \\
Scénarios Behat des modules syllabus & 6 scénarios, 102 étapes \\
Style Moodle (\code{phpcs}), PHPDoc, validation, points de mise à jour (CI) & sans remarque \\
Scénario de bout en bout des modules syllabus, code corrigé & 51/51 contrôles \\
Balayage des droits : 59 pages $\times$ 3 comptes & 177 appels, 27 ouvertures attendues \\
\bottomrule
\end{tabularx}
\end{center}

\paragraph{Balayage des droits.} Chaque page des cinq modules est appelée avec l'identifiant
d'une activité par trois comptes : un étudiant d'une autre promotion (non inscrit), un étudiant
inscrit, le responsable d'un autre module. Les 27 pages qui ne refusent pas sont toutes
légitimes : pages de l'étudiant (vue, catalogue, déclaration, son supplément, ses éléments, sa
fiche de module en PDF, « Mes modules »), vérification publique d'un supplément par son code,
pages des enseignants du cours (tableau de pilotage, validation, propositions, matrice et
progression ouvertes aux enseignants), page des compétences d'un module en lecture seule. Le
balayage ne fait que des requêtes de consultation : les actions sont couvertes par les tests
unitaires et la revue du code.

\section{Constats corrigés}

\subsection{Sécurité}

\begin{xltabular}{\linewidth}{@{}l X X@{}}
\toprule
\textbf{Gravité} & \textbf{Constat} & \textbf{Correction} \\
\midrule
\endhead
\grav{moyenne} & \code{supplement.php} et \code{dashboard.php} : l'identifiant d'étudiant de
l'adresse n'était pas contrôlé ; un détenteur de \code{mod/ep:viewall} lisait le supplément
(PDF avec date de naissance) ou la situation de tout compte du site. & Contrôle de l'inscription
à l'activité, comme à l'émission. Banc : PDF obtenu avant, « Invalid user ID » après. \\
\grav{moyenne} & Déclaration en attente : le type changeait, pas l'étape en attente ; le
référent tranchait un type que seule la DEVE valide. & Étape recalculée sur le circuit du
nouveau type (test \texttt{test\_type\_change\_\-follows\_new\_flow}). \\
\grav{haute} & Avis : un responsable fixait librement la fenêtre ; il pouvait la fermer pour
lire les résultats, la rouvrir, ou ouvrir dans une autre année pour archiver le recueil en
cours. & Dès le premier avis, seule la DEVE déplace la fenêtre, par la page du recueil comme par
les paramètres de l'activité (test \texttt{test\_window\_locked\_\-after\_first\_response}). \\
\grav{haute} & Avis : l'événement d'envoi gardait l'auteur et l'heure ; rapproché de
l'évolution des résultats, il désignait l'auteur de chaque avis. & Événement anonyme : le
journal ne le montre qu'aux administrateurs (\code{moodle/site:viewanonymousevents}). \\
\grav{moyenne} & Résultats comparés du suivi : seul le droit du suivi était vérifié, pas celui
de chaque module. & Modules filtrés sur \code{mod/syllabus:manage}, comme la liste et la
relecture. \\
\grav{basse} & Libellés des listes déroulantes (nom et numéro d'étudiant, code de compétence,
période) non échappés : \code{html_writer::select()} ne le fait pas. & \code{s()} sur ces
libellés. \\
\grav{basse} & Exports CSV : une cellule commençant par \code{=}, \code{+}, \code{-} ou
\code{@} s'exécutait comme formule dans un tableur ; le cœur ne le neutralise que dans ses
versions récentes. & Apostrophe ajoutée à ces cellules dans les 15 lignes d'export des deux
modules. \\
\grav{moyenne} & Suppressions de questions, types, référentiels et compétences confirmées en
JavaScript seulement : sans JavaScript, un clic supprimait. & Page de confirmation du serveur
comme dans \code{mod_ep} (\texttt{syllabus\_require\_\-confirmation()}). \\
\bottomrule
\end{xltabular}

\subsection{Intégrité des données}

\begin{xltabular}{\linewidth}{@{}l X X@{}}
\toprule
\textbf{Gravité} & \textbf{Constat} & \textbf{Correction} \\
\midrule
\endhead
\grav{haute} & Restauration : le sceau des suppléments n'était pas sauvegardé ; après
suppression et restauration d'une activité (corbeille), les diplômes remis étaient déclarés
« non reconnus ». & Sceau sauvegardé ; sur le même site, s'il est authentique et que le numéro
d'origine est libre, il est recalculé ; toute autre copie reste non reconnue (deux tests). \\
\grav{moyenne} & Restauration : les décisions d'un crédit écarté étaient rattachées au crédit
restauré juste avant. & \code{SKIP_ALL_CHILDREN}. \\
\grav{moyenne} & Rentrée appliquée deux fois depuis deux onglets ou par un envoi répété (année
N vers N+2). & Mise à jour conditionnée à l'année de départ confirmée (test). \\
\grav{moyenne} & Transfert d'un étudiant : ses crédits de stage étaient supprimés comme
orphelins par la synchronisation de la cible. & Convertis en saisies DEVE si la cible ne lit pas
les mêmes stages, laissés à sa synchronisation sinon (test). \\
\grav{moyenne} & Suppression d'un EP du catalogue : versions du syllabus, séances et
prérequis restaient orphelins ; le contrôle des inscriptions n'était pas sous verrou ; le bouton
échouait toujours (clé de session absente du lien). & \code{ep_delete_activity()} sous verrou et
en transaction ; page de confirmation (test). \\
\grav{moyenne} & Émission groupée des suppléments : une erreur faisait relancer toute la tâche,
qui réémettait les suppléments déjà émis (nouveaux numéros). & Un étudiant en échec n'arrête pas
les autres ; ceux dont le supplément a été émis depuis la demande sont sautés (test). \\
\grav{basse} & Copie d'un cours sans données utilisateurs : la fenêtre des avis gardait les
anciennes dates ; les étudiants du nouveau cours étaient prévenus d'un recueil passé. & Recueil
sans fenêtre après une telle copie (test). \\
\grav{basse} & Tâche d'ouverture des avis : un arrêt en cours d'envoi renvoyait tous les
messages du module. & Module marqué avant l'envoi. \\
\grav{basse} & Transfert : le second verrou était pris hors du bloc qui libère le premier. &
Pris dans le bloc. \\
\bottomrule
\end{xltabular}

\subsection{Performance et portabilité}

\begin{xltabular}{\linewidth}{@{}l X X@{}}
\toprule
\textbf{Gravité} & \textbf{Constat} & \textbf{Correction} \\
\midrule
\endhead
\grav{haute} & Intitulé des crédits de stage sans troncature (thème et structure jusqu'à
255 caractères chacun) : l'insertion échouait, et avec elle chaque page qui synchronise. & Borné
à 255 caractères (test) ; rédigé dans la langue du site, pour qu'une consultation dans une autre
langue ne réécrive pas tous les crédits. \\
\grav{moyenne} & Décisions en lot : les courriels partaient pendant que le verrou de l'activité
était pris (20 à 60\,s pour une promotion). & Messages gardés puis envoyés après validation et
libération du verrou ; rien n'est envoyé si le lot échoue (test). \\
\grav{moyenne} & Tâche d'ouverture des avis : une requête par étudiant. & Une requête par
module. \\
\grav{moyenne} & \code{AVG} sur un entier : moyenne tronquée sous SQL Server. & Moyenne
décimale. \\
\grav{basse} & \code{DISTINCT} sur la fiche utilisateur (champs texte, refusé par Oracle) ;
\code{<> ''} (Oracle) ; alias \code{uid} (réservé sous Oracle). & Sous-requête,
\code{sql_isnotempty()}, alias \code{userid}. \\
\bottomrule
\end{xltabular}

\subsection{Langues, accessibilité, documentation}

\begin{itemize}
  \item Listes de filtres : le nom lu par les lecteurs d'écran était le nom technique du
    paramètre (« status ») ; c'est désormais le libellé de la liste.
  \item Cases à cocher des responsables d'un EP et des activités liées à une synthèse nommées ;
    indicateurs d'étape (coche, cercle) doublés d'un texte « Fait » / « À faire ».
  \item Tableaux larges (gabarit commun de \code{mod_ep}, liste des modules, résultats,
    catalogue, synthèse) dans un conteneur défilant : la page ne défile plus en largeur sur
    téléphone.
  \item Chaînes absentes du cœur 4.5 (\code{nonewmodules}, \code{cannotcreatezip}) remplacées :
    elles s'affichaient entre crochets.
  \item README et INSTALL au vocabulaire « Avis des étudiants » ; règle de la fenêtre et
    anonymat du journal documentés ; guides au numéro des modules (guide syllabus 1.7).
  \item \code{\$plugin->supported = [403, 405]} dans les cinq modules ; CHANGELOG corrigé sur
    la datation des versions et complété ; trois fichiers de tests rendus exécutables seuls.
\end{itemize}

\section{Constats restants et recommandations}\label{sec:residuel}

\subsection{Recommandations prioritaires}

\begin{description}
  \item[CI : mise à jour depuis la version en production] (\grav{haute}). La CI n'installe qu'à
    neuf ; la conversion des données des sites en 0.1.0 n'a été vérifiée que sur le banc
    (\code{migration_avant.php}, \code{migration_apres.php}). Ajouter un job qui installe
    \code{main}, crée des données, passe au code de la branche, lance la mise à jour et rejoue
    les contrôles.
  \item[Behat de \code{mod_ep}] (\grav{haute}). Aucun scénario pour les 43 pages de
    \code{mod_ep}, alors que le CHANGELOG 0.2.0 en annonce. Le balayage des droits du banc
    couvre la consultation ; trois scénarios sont recommandés : justificatif d'un autre étudiant
    refusé, supplément émis visible du seul étudiant et de la DEVE, pages DEVE refusées à un
    enseignant.
  \item[Réglages communs des avis] (\grav{moyenne}). Seuil d'anonymat, résultats des
    enseignants, relecture et questions communes valent pour tout le site, mais se règlent
    depuis n'importe quelle activité de suivi. C'est le choix demandé (tout depuis le suivi,
    sans l'administration du site) ; il suppose de ne donner \code{mod/syllabusstaff:manageall}
    qu'à la DEVE et de garder un seuil d'au moins 3. Une capacité au niveau du site est
    possible si plusieurs services gèrent des catégories distinctes.
  \item[Anonymat en base] (\grav{moyenne}, risque résiduel). Participation et réponse sont
    insérées l'une après l'autre : leurs identifiants séquentiels se correspondent, en base et
    dans une sauvegarde avec données utilisateurs. Aucune page ni aucun export ne les relie ; le
    risque est documenté dans l'INSTALL : accès à la base et aux sauvegardes restreint. Découpler
    l'ordre (réponses insérées par lots mélangés) demanderait une file d'attente.
\end{description}

\subsection{Autres constats, sans correction}

\begin{xltabular}{\linewidth}{@{}l X X@{}}
\toprule
\textbf{Gravité} & \textbf{Constat} & \textbf{Raison, recommandation} \\
\midrule
\endhead
\grav{moyenne} & Restauration d'une archive globale dans une instance liée aux mêmes stages :
ECTS de stage comptés deux fois (plausible). & Cas rare (remise à zéro puis restauration) ;
écarter les crédits de stage dans ce cas. \\
\grav{moyenne} & Oracle : colonnes \code{comment}, \code{number}, \code{level}, mots réservés. &
Oracle n'est ni visé ni testé ; les renommer est un changement de schéma, à faire seulement si
Oracle devient une cible. \\
\grav{basse} & Responsables : tout responsable désigne ou retire des co-responsables ; dans un
module sans responsable, un enseignant peut se désigner. & Conforme au guide ; la DEVE voit les
responsables dans la liste des modules. \\
\grav{basse} & Progression, livret : modules de cours masqués visibles aux enseignants du
suivi. & Filtrer sur la visibilité du cours si des cours en préparation doivent rester
discrets. \\
\grav{basse} & Effacement RGPD d'un étudiant : il peut redonner son avis la même année. &
Effacement rare et en fin de parcours ; documenté. \\
\grav{basse} & Imports : responsables et validateurs reconnus sur tout le site. & Aucun gain de
droits ; limiter aux enseignants du cours si besoin. \\
\grav{basse} & Concurrences sans corruption : double décision d'année (erreur d'index unique),
deux brouillons de syllabus de même numéro, création simultanée du secret du sceau. & Cas
improbables ; verrouiller si observés. \\
\grav{basse} & Relances fondées sur la date de création de la demande ; DEVE non inscrite au
cours jamais relancée. & Inscrire la DEVE au cours (recommandé par l'INSTALL). \\
\grav{basse} & Réglage du premier mois de l'année changé en cours d'année : le recueil suivant
archive le précédent. & Ne changer ce réglage qu'entre deux années ; documenté. \\
\grav{basse} & Liens de compétences des cours sans module absents des sauvegardes de cours ;
périmètre de catégorie restauré sur un autre site. & À refaire après une copie d'un site à
l'autre. \\
\grav{basse} & Réinitialisation : empreinte de synchronisation des stages gardée (crédits de
stage recalculés au plus tard à la tâche de nuit) ; suppression par sous-requête lente sous
MariaDB pour de très gros volumes. & Sans conséquence fonctionnelle. \\
\grav{basse} & Livrets : deux premiers téléchargements simultanés produisent tous deux le PDF ;
renommage d'un responsable visible le lendemain ; archive ZIP des suppléments lue fichier par
fichier. & Coût mesuré acceptable (rapports de charge). \\
\grav{basse} & Messages rédigés dans la langue de l'expéditeur ou du site, non du destinataire.
& Site monolingue ; à revoir pour un site bilingue. \\
\grav{basse} & Accessibilité : pas de légende ni d'en-tête de ligne dans les tableaux, sens de
tri sans \code{aria-sort}, niveaux de titres sautés, actions de ligne sans contexte. & Lisible
au lecteur d'écran ; à reprendre avec une revue RGAA. \\
\grav{basse} & Langues : concaténations avec « : » codé en dur, « DEVE » et « academic office »
mêlés en anglais, descriptions d'événements de \code{mod_ep} et traces des tâches en français,
19 chaînes inutilisées. & Sans effet en français ; à harmoniser avec la traduction anglaise. \\
\grav{basse} & Retrait de fenêtre en lot confirmé en JavaScript seulement. & Retrait réversible
(la fenêtre se refixe), réponses conservées. \\
\grav{moyenne} & CI : dépôt StageCompagnon pris sur \code{main} sans épinglage, outils non
épinglés, ni PHP~8.3, ni MySQL, ni SQL Server ; seuil de couverture absent. & Épingler un
commit de StageCompagnon et ajouter une ligne PHP~8.3 lors d'une prochaine révision de la CI. \\
\grav{basse} & Tests : quelques assertions faibles (exception générique, crédit préexistant non
vérifié) ; test annuel sensible à la minute qui suit le 1\textsuperscript{er} septembre. &
Probabilité très faible ; à resserrer au fil de l'eau. \\
\grav{basse} & Banc : prérequis (Python \code{requests}, Playwright, \code{pdftoppm}) et en-têtes
GPL des scripts ; traduction française minimale des guides sans attribution. & Prérequis
d'exécution précisés dans le README du banc ; le reste à compléter. \\
\bottomrule
\end{xltabular}

\section{Conclusion}

Les contrôles de base sont solides et uniformes : chaque page exige la connexion au cours et un
droit dans le bon contexte, toutes les écritures passent par la clé de session, tout le SQL est
paramétré, l'échappement est cohérent, aucun service web ni code JavaScript n'est exposé. Les
défauts trouvés tenaient à des chemins secondaires -- un identifiant de la requête non relié à
l'activité, un champ non recalculé, une donnée non sauvegardée -- et à l'anonymat des avis face
aux traces indirectes (journal, ouverture et fermeture de la fenêtre).

Les 28 corrections sont couvertes par 7 tests unitaires de plus (190 au total), le balayage
des droits et le scénario de bout en bout rejoués sur le banc. La branche peut être fusionnée ;
les deux recommandations de CI (mise à jour depuis la version en production, Behat de
\code{mod_ep}) sont à traiter ensuite, avant la prochaine évolution du schéma.

\end{document}
